\section{IDT 门描述符的 8 字节布局}

每个 IDT 条目（gate descriptor）是 8 字节（64 位），格式和 GDT 描述符一样"碎片化"：

\subsection{内存布局}

\begin{figure}[H]
  \centering
  % figures/ch04/fig01-idt-gate-layout.tex — auto-extracted from chapter source
\begin{tikzpicture}[node distance=0cm]
  \node[memcell, minimum width=6cm, fill=green!10] (b0) {offset\_low [7:0]};
  \node[memaddr, left=0.3cm of b0] {+0};

  \node[memcell, minimum width=6cm, fill=green!10, above=0cm of b0] (b1) {offset\_low [15:8]};
  \node[memaddr, left=0.3cm of b1] {+1};

  \node[memcell, minimum width=6cm, fill=blue!10, above=0cm of b1] (b2) {selector [7:0]};
  \node[memaddr, left=0.3cm of b2] {+2};

  \node[memcell, minimum width=6cm, fill=blue!10, above=0cm of b2] (b3) {selector [15:8]};
  \node[memaddr, left=0.3cm of b3] {+3};

  \node[memcell, minimum width=6cm, fill=gray!15, above=0cm of b3] (b4) {zero (must be 0)};
  \node[memaddr, left=0.3cm of b4] {+4};

  \node[memcell, minimum width=6cm, fill=red!10, above=0cm of b4] (b5) {type\_attr};
  \node[memaddr, left=0.3cm of b5] {+5};

  \node[memcell, minimum width=6cm, fill=green!10, above=0cm of b5] (b6) {offset\_high [23:16]};
  \node[memaddr, left=0.3cm of b6] {+6};

  \node[memcell, minimum width=6cm, fill=green!10, above=0cm of b6] (b7) {offset\_high [31:24]};
  \node[memaddr, left=0.3cm of b7] {+7};

  % Braces
  \draw[decorate, decoration={brace, amplitude=8pt}]
    ([xshift=0.3cm]b1.north east) -- ([xshift=0.3cm]b0.south east)
    node[midway, right=10pt, font=\small, text=green!50!black] {handler 地址低 16 位};
  \draw[decorate, decoration={brace, amplitude=8pt}]
    ([xshift=0.3cm]b3.north east) -- ([xshift=0.3cm]b2.south east)
    node[midway, right=10pt, font=\small, text=blue!50!black] {代码段 selector};
  \draw[decorate, decoration={brace, amplitude=8pt}]
    ([xshift=0.3cm]b7.north east) -- ([xshift=0.3cm]b6.south east)
    node[midway, right=10pt, font=\small, text=green!50!black] {handler 地址高 16 位};
\end{tikzpicture}

  \caption{IDT 门描述符的 8 字节内存布局}
  \label{fig:idt-gate-layout}
\end{figure}

\subsection{C 结构体}

\codefile{src/include/arch/idt.h}
\begin{lstlisting}[style=cstyle]
typedef struct __attribute__((packed)) {
    uint16_t offset_low;    // handler 地址低 16 位 [0:15]
    uint16_t selector;      // 代码段选择子（0x08 = 内核代码段）
    uint8_t  zero;          // 保留字段，必须为 0
    uint8_t  type_attr;     // P(1) | DPL(2) | 0(1) | Type(4)
    uint16_t offset_high;   // handler 地址高 16 位 [16:31]
} idt_entry_t;
\end{lstlisting}

\begin{thinkbox}
handler 地址为什么也被拆成 low 和 high 两段？和 GDT 描述符一样，
这是 286→386 向后兼容的历史包袱。286 的 IDT 门只有 16 位 offset。
\end{thinkbox}

\subsection{\texttt{type\_attr} 字节逐 bit 详解}

\begin{figure}[H]
  \centering
  % figures/ch04/fig02-type_attr-8-bit.tex — auto-extracted from chapter source
\begin{tikzpicture}
  \foreach \i/\label in {0/Type$_0$, 1/Type$_1$, 2/Type$_2$, 3/Type$_3$,
                          4/0, 5/DPL$_0$, 6/DPL$_1$, 7/P} {
    \node[bitcell, fill=white] (t\i) at (\i*1.0, 0) {\label};
    \node[font=\tiny\ttfamily, above=0.1cm of t\i] {bit \i};
  }
  \draw[decorate, decoration={brace, amplitude=5pt, mirror}]
    (0, -0.4) -- (3.5, -0.4) node[midway, below=5pt, font=\small] {Type (4 bit)};
  \draw[decorate, decoration={brace, amplitude=5pt, mirror}]
    (4.5, -0.4) -- (6.5, -0.4) node[midway, below=5pt, font=\small] {DPL};
\end{tikzpicture}

  \caption{type\_attr 字节的 8 个 bit}
\end{figure}

\begin{bitfield}
  \bit{7}{P (Present) = 1：此门有效。P=0 时触发该向量会引发 \#NP。}
  \bit{5--6}{DPL (Descriptor Privilege Level)：
    允许通过 \texttt{int N} 软件触发此门的最低特权级。
    DPL=0：只有 Ring 0 能用 \texttt{int N} 触发。
    DPL=3：Ring 3（用户态）也能触发（如 \texttt{int 0x80} 系统调用）。
    \textbf{注意}：硬件异常不受 DPL 限制。CPU 除零时无论 DPL 都会进入 handler。}
  \bit{4}{固定为 0（区分系统段描述符和门描述符）。}
  \bit{0--3}{Type：
    \begin{itemize}
      \item \texttt{1110b} = \hex{E} = 32-bit Interrupt Gate：进入时\textbf{自动清 IF}（关中断）
      \item \texttt{1111b} = \hex{F} = 32-bit Trap Gate：进入时\textbf{不清 IF}
    \end{itemize}}
\end{bitfield}

我们使用 \hex{8E}：

$$\hex{8E} = 1\underbrace{00}_{DPL=0}\underbrace{0}_{固定}\underbrace{1110}_{Interrupt Gate}$$

\begin{designbox}
Interrupt Gate vs Trap Gate——唯一区别是进入 handler 时是否自动清 \reg{EFLAGS}.IF。

\textbf{Interrupt Gate}（\hex{8E}）：自动关中断。适合异常处理——
处理 \#PF 时我们不希望被键盘中断打断，否则嵌套的中断可能在异常尚未修复时
访问出错的页面。

\textbf{Trap Gate}（\hex{8F}）：不关中断。将来 syscall gate（\texttt{int 0x80}）
可以考虑用 Trap Gate，让系统调用期间仍能响应硬件中断。
\end{designbox}

\subsection{\texttt{idt\_set\_gate}：填充一个门描述符}

\codefile{src/kernel/arch/idt.c}
\begin{lstlisting}[style=cstyle]
static void idt_set_gate(uint8_t vector, void (*handler)(void),
                         uint16_t selector, uint8_t type_attr) {
    uint32_t addr = (uint32_t)(uintptr_t)handler;

    idt[vector].offset_low  = (uint16_t)(addr & 0xFFFF);
    //   addr = 0xC0101234
    //   addr & 0xFFFF = 0x1234 → offset_low

    idt[vector].selector    = selector;
    //   0x08 (内核代码段)

    idt[vector].zero        = 0;
    //   保留字段，必须为 0

    idt[vector].type_attr   = type_attr;
    //   0x8E = Interrupt Gate, DPL=0, Present

    idt[vector].offset_high = (uint16_t)((addr >> 16) & 0xFFFF);
    //   (addr >> 16) & 0xFFFF = 0xC010 → offset_high
}
\end{lstlisting}

CPU 触发向量 $N$ 时，从 IDT[$N$] 拼出 handler 地址：

$$\text{handler} = (\text{offset\_high} \ll 16) \;|\; \text{offset\_low}$$

然后用 selector（\hex{08}）去 GDT 查代码段描述符，切换到内核代码段执行 handler。

\subsection{IDTR 结构体}

和 GDTR 格式完全相同：6 字节，limit + base。

\codefile{src/include/arch/idt.h}
\begin{lstlisting}[style=cstyle]
typedef struct __attribute__((packed)) {
    uint16_t limit;  // IDT 总字节数 - 1
    uint32_t base;   // IDT 数组的线性地址
} idtr_t;
\end{lstlisting}

256 个门 $\times$ 8 字节 = 2048 字节，所以 \texttt{limit = 2047}。

% ============================================================================

